% ED26-SCI-GelmanMengStern1996-001; original main article, pp. 733--760.
\chapter{Posterior Predictive Model Criticism}
\label{sec:p15-posterior-predictive-model-criticism}

A predictive check asks whether a specified aspect of the observations is
typical of replications generated by the fitted model. Its inputs are a
sampling model, a prior, the observed data, a precise replication scheme,
and a discrepancy chosen for the inferential purpose. Its output is a
reference distribution, together with a comparison of observed and replicated
discrepancies. A favorable comparison leaves unexamined aspects of the model
unexamined. It supplies neither posterior odds between models nor a probability
that the model is true.

Gelman, Meng and Stern extend posterior predictive checks from statistics
depending only on observations to discrepancies depending on observations
and unknown parameters \cite[secs.~1--4]{GelmanMengStern1996Main}.
The following treatment develops that construction, its computational and
calibration qualifications, the article's theoretical examples and its three
applications. The probability foundations are
\cref{thm:p15-bayes-posterior-kernel,thm:p15-bayes-predictive-kernel,thm:p15-conditional-variance};
model odds have their separate definition in
\cref{thm:p15-model-posterior-odds}.

\section{The model, the replication and the discrepancy}
\label{sec:p15-ppc-construction}

\begin{hbdefinition}{def:p15-ppc-model}{Observed and replicated data}
Work on standard Borel parameter and data spaces \(\Theta\) and
\(\mathcal Y\). The model \(H\) supplies a sampling kernel \(P_\theta\),
a prior \(\Pi\), and a version \(\Pi_y\) of the posterior at the observed
data \(y\). The posterior must be a probability measure. An improper prior
may be used only when its likelihood-weighted measure has a finite, positive
normalizing constant; prior predictive claims below additionally require
a proper prior.

A replication kernel \(R_a(\theta,dy')\) describes the experiment to be
repeated, where \(a=A(y)\) is the observed value of explicitly fixed auxiliary
information. For example, \(a\) may specify sample size, exposure times,
covariates or a conditioned total. For a conditional replication, use a
specified version of
\[
 R_a(\theta,B)=P_\theta(Y'\in B\mid A(Y')=a).
\]
Conditioning on a continuous auxiliary variable requires a regular conditional
kernel; it is not division by the probability of a null event. The posterior
\(\Pi_y\) uses the original observation likelihood. Choosing a conditional
replication does not by itself authorize replacing that likelihood.
With no additional conditioning, \(R_a=P_\theta\).
\index{posterior predictive checking!replication}
\end{hbdefinition}

\begin{hbdefinition}{def:p15-ppc-discrepancy}{Realized discrepancy and reference probability}
A statistic is a measurable map \(T:\mathcal Y\to\mathbb R\).
A discrepancy is a jointly measurable map
\(D:\mathcal Y\times\Theta\to\mathbb R\). Its realized value
\(D(y,\theta)\) can remain unknown after observing \(y\).
In a posterior replication, the same parameter value is used in both
\(D(y,\theta)\) and \(D(y',\theta)\). Define
\[
 q_D(y,\theta)=\int
  \mathbf1\{D(y',\theta)\geq D(y,\theta)\}\,R_{A(y)}(\theta,dy'),
 \qquad p_B(y)=\int q_D(y,\theta)\,\Pi_y(d\theta).
\]
The convention is an upper tail including ties. A lower tail is obtained by
reversing the comparison. A two-sided assessment requires a specified ordering
of discrepancies; it is not determined by the words ``two-sided'' alone.
\index{realized discrepancy}\index{posterior predictive p-value}
\end{hbdefinition}

\begin{hbtheorem}{thm:p15-ppc-joint}{proved=annex:p15.ppc-joint}{The joint predictive reference construction}
The kernel
\[
 Q_y(d\theta,dy')=\Pi_y(d\theta)R_{A(y)}(\theta,dy')
\]
is a probability law. Its data marginal is
\(M_y(B)=\int R_{A(y)}(\theta,B)\Pi_y(d\theta)\), and
\[
 p_B(y)=Q_y\{(\theta,y'):D(y',\theta)\geq D(y,\theta)\}\in[0,1].
\]
For \(D(y,\theta)=T(y)\), this is the upper tail of the predictive
distribution of \(T(Y')\), evaluated at \(T(y)\).
For a parameter-dependent discrepancy it is a probability about a pair;
it is generally not the replicated discrepancy's tail above one fixed observed
number. Replacing the shared parameter by independent posterior draws
in the two discrepancies changes the comparison law.
\end{hbtheorem}

\begin{hbdefinition}{def:p15-ppc-reduced-discrepancies}{Minimum and posterior-average statistics}
When the operations exist and give measurable functions, set
\[
 D_{\min}(y)=\inf_{\theta\in\Theta}D(y,\theta),\qquad
 D_{\rm avg}(y)=\int D(y,\theta)\Pi_y(d\theta).
\]
Use a minimum only when the infimum is attained; require integrability for
the average. Both are statistics, so their reference values at \(y'\)
are \(D_{\min}(y')\) and \(D_{\rm avg}(y')\).
The latter integrates against \(\Pi_{y'}\), the posterior refitted to
the replicate, rather than \(\Pi_y\). Consequently, an optimization or a
new posterior calculation can be required for every replicated data set.
\end{hbdefinition}

\section{Computation and graphical comparison}
\label{sec:p15-ppc-computation}

\begin{hbmethod}{def:p15-ppc-simulation}{Posterior simulation with paired discrepancies}
For each retained posterior draw \(\theta_j\), generate
\(y'_j\sim R_{A(y)}(\theta_j,\cdot)\). Record the ordered pair
\(\bigl(D(y,\theta_j),D(y'_j,\theta_j)\bigr)\).
Plot observed discrepancy horizontally and replicated discrepancy vertically,
using the same scale, with the diagonal as the comparison boundary.
The fraction on or above the diagonal estimates \(p_B(y)\).
Also examine replicated data on their original scale: a small residual tail
probability does not identify which observations or features are responsible.
The same draws may be evaluated with several discrepancies, with their choices
and purposes recorded. If \(q_D(y,\theta_j)\) is available analytically,
its average removes the extra replication noise.
\end{hbmethod}

\begin{hbtheorem}{thm:p15-ppc-simulation}{proved=annex:p15.ppc-simulation}{Monte Carlo targets and the benefit of conditional averaging}
Suppose \((\theta_j,y'_j)\), \(1\leq j\leq J\), are independent draws
from \(Q_y\), with \(y\) fixed. Write
\(I_j=\mathbf1\{D(y'_j,\theta_j)\geq D(y,\theta_j)\}\),
\(\widehat p=J^{-1}\sum_j I_j\), and
\(\widetilde p=J^{-1}\sum_jq_D(y,\theta_j)\). Then
\[
 E(\widehat p\mid y)=E(\widetilde p\mid y)=p_B(y),\quad
 \operatorname{Var}(\widehat p\mid y)=\frac{p_B(y)(1-p_B(y))}{J},
\]
\[
 \operatorname{Var}(\widetilde p\mid y)
 \leq\operatorname{Var}(\widehat p\mid y),\qquad
 \Pr(|\widehat p-p_B|\geq\epsilon\mid y)
 \leq\frac1{4J\epsilon^2}\quad(\epsilon>0).
\]
These variance formulas require independence. For stationary, correlated
indicators with covariance \(\gamma_h=\operatorname{Cov}(I_j,I_{j+h}\mid y)\),
\[
 \operatorname{Var}(\widehat p\mid y)
 =J^{-2}\left[J\gamma_0+2\sum_{h=1}^{J-1}(J-h)\gamma_h\right].
\]
An approximate posterior sampler introduces approximation and convergence
questions beyond this Monte Carlo calculation.
\end{hbtheorem}

The article permits dependent posterior simulations, but a binomial standard
error cannot be assigned to them solely because each plotted point is a pair.
Failure to explore a mixture's separated modes can bias the reference
distribution. Replacing a difficult model fit by a local optimum similarly
changes a computed minimum-discrepancy check. Simulation diagnostics give
evidence about these computations, not a theorem that finite chains have mixed.

\section{Calibration and an exactly uninformative mean check}
\label{sec:p15-ppc-calibration}

\begin{hbtheorem}{thm:p15-ppc-pivot}{proved=annex:p15.ppc-pivot}{A pivot and a posterior conditional expectation}
If the conditional law of \(T(Y')\) under \(R_a(\theta,\cdot)\) is
the same law \(F_a\) for every posterior-relevant \(\theta\), then the
predictive law is \(F_a\) for every prior yielding that posterior.
If an observed \(Y\), under the sampling model being assessed conditional
on \(A(Y)=a\), also has continuous statistic distribution \(F_a\),
the tail \(1-F_a(T(Y))\) is uniform on \([0,1]\).

For a different statement, let the random parameter \(\vartheta\) have
proper prior \(\Pi\), take no auxiliary conditioning, and let
\(Y,Y'\) be conditionally independent with law \(P_\vartheta\). Assume
\(D(Y,\theta)\) has a continuous distribution under each \(P_\theta\).
Then, under this joint prior model,
\[
 E[p_B(Y)]=\tfrac12,\qquad
 \operatorname{Var}(p_B(Y))\leq\tfrac1{12}.
\]
These are prior-model statements, not a uniform distribution theorem for
\(p_B\), and not fixed-parameter type-I-error guarantees.
\end{hbtheorem}

\begin{hbtheorem}{thm:p15-ppc-normal-mean}{proved=annex:p15.ppc-normal-mean}{Nonuniformity in a normal model}
Let \(Y_i\mid\theta\) be independent \(N(\theta,\sigma^2)\), with
known \(\sigma>0\), \(n\geq1\), prior \(N(0,\tau^2)\), and
\(\tau>0\). Put \(s^2=\sigma^2/n\),
\(w=\tau^2/(\tau^2+s^2)\), and \(v=\tau^2s^2/(\tau^2+s^2)\).
For the statistic \(T(y)=\bar y\) and a replicate of the same size,
\[
 \theta\mid y\sim N(w\bar y,v),\qquad
 p_B(y)=1-\Phi\left(\frac{(1-w)\bar y}{\sqrt{s^2+v}}\right),
\]
where \(\Phi\) is the standard normal distribution function.
Under the joint prior model, the argument of \(\Phi\) has distribution
\(N(0,s^2/(s^2+2\tau^2))\), so \(p_B(Y)\) is not uniform.
With the improper flat prior \(d\theta\), the posterior is proper,
\(\theta\mid y\sim N(\bar y,s^2)\), and \(p_B(y)=1/2\) for every
data vector. This flat-prior mean check cannot diagnose any feature of
the data, including extreme dispersion or dependence.
\end{hbtheorem}

The last conclusion concerns this particular statistic, not all predictive
checks. For example, residual extremes, unequal subgroup variation and serial
patterns ask different questions. Using data in the posterior and in the
discrepancy is part of the stated construction; it also means that classical
uniform calibration cannot be silently imported.

\section{Squared residuals and fitted linear models}
\label{sec:p15-ppc-linear}

For independent observations with finite means \(\mu_i(\theta)\) and
strictly positive variances \(v_i(\theta)\), define the Pearson discrepancy
\[
 X^2(y,\theta)=\sum_{i=1}^n
       \frac{(y_i-\mu_i(\theta))^2}{v_i(\theta)}.
\]
Its expectation at the generating parameter is \(n\), but its name does
not imply a chi-square distribution. Under independent normal errors with
known variances it is exactly chi-square; for other sampling laws any
chi-square approximation needs its own justification. Sparse Poisson counts,
constraints and mixtures are especially poor grounds for assigning degrees
of freedom merely by counting displayed parameters.

Under the article's assumption that the reference law of
\(X^2(Y',\theta)\) is approximately \(\chi^2_n\) at relevant
parameter values, its reference calculation is
\[
 p_B(y)\approx\int\Pr\{\chi^2_n\geq X^2(y,\theta)\}\Pi_y(d\theta).
\]
The approximation must remain accurate over the posterior-supported
parameters. For a diffuse prior in a regular linear approximation, the
additional posterior discrepancy above the minimum is approximated by
\(\chi^2_k\), yielding an average shift of \(k\) and a difference
reference involving independent chi-squares. The article uses these as
diagnostic approximations, without a uniform error bound. The next result
states and proves a setting where every step is exact.

\begin{hbtheorem}{thm:p15-ppc-linear}{proved=annex:p15.ppc-linear}{An exact Gaussian form of the three discrepancy calculations}
Let \(Y\mid\beta\sim N_n(B\beta,I_n)\), where the real \(n\times k\)
matrix \(B\) has full column rank and \(1\leq k<n\). Give \(\beta\)
the improper flat prior on \(\mathbb R^k\). Write
\(P=B(B^TB)^{-1}B^T\), \(r(y)=\|(I_n-P)y\|^2\), and
\(D(y,\beta)=\|y-B\beta\|^2\). Then
\[
 D_{\min}(y)=r(y),\qquad D_{\rm avg}(y)=r(y)+k.
\]
The predictive tails based on either of these statistics are
\(\Pr(\chi^2_{n-k}\geq r(y))\). The realized-discrepancy tail is
\[
 p_B(y)=\Pr(U-V\geq r(y)),\qquad
 U\sim\chi^2_n,\quad V\sim\chi^2_k,\quad U\ \hbox{independent of }V.
\]
The two reference distributions have the same mean \(n-k\), with variances
\(2(n-k)\) and \(2(n+k)\), respectively. There is no universal ordering
of their upper tails. In particular, at \(r(y)=0\), the minimum tail is
one, whereas \(\Pr(U-V\geq0)<1\).
\end{hbtheorem}

The article's illustration takes \(n=250\), \(k=200\), and
\(r(y)=80\). Relative to \(\chi^2_{50}\), 80 is three standard
deviations above the mean 50. Relative to the difference of independent
\(\chi^2_{250}\) and \(\chi^2_{200}\), it is one standard deviation
above that same mean. The checks measure distance from the best fitting
mean and distance from posterior-supported means differently. This contrast
does not prove that a discrepancy harmless for one use is harmless for another.
Numerical integration of the exact Gaussian laws gives upper tails
approximately \(0.00448266\) and \(0.15813111\), respectively.
With a different prior, the minimum residual remains pivotal in this Gaussian
model, while realized and posterior-average discrepancies can change.

\section{A count example: a correct total can hide dispersion}
\label{sec:p15-ppc-counts}

\begin{hbtheorem}{thm:p15-ppc-poisson}{proved=annex:p15.ppc-poisson}{Predictive totals and a conditional dispersion check}
Let \(n\geq1\) be an integer and let \(Y_1,\ldots,Y_n\mid\lambda\) be independent Poisson variables
with mean \(\lambda>0\), and take a gamma prior with shape \(a>0\)
and rate \(b>0\), whose density is proportional to
\(\lambda^{a-1}e^{-b\lambda}\). Write \(S=\sum_i y_i\).
The posterior is gamma with shape \(a+S\) and rate \(b+n\). Under
an unconditioned replicate of the same size, the total \(S'\) has law
\[
 \Pr(S'=r\mid y)=
 \frac{\Gamma(a+S+r)}{\Gamma(a+S)r!}
 \left(\frac{b+n}{b+2n}\right)^{a+S}
 \left(\frac n{b+2n}\right)^r,
 \qquad r=0,1,\ldots.
\]
Thus every total or mean check based on this law gives identical results
for all vectors with the same \(S\).

If instead the replication holds the total \(S>0\) fixed, then
\(Y'\mid S\sim\operatorname{Multinomial}(S;1/n,\ldots,1/n)\),
independently of \(\lambda\). For
\[
 T(y)=\sum_{i=1}^n\frac{(y_i-S/n)^2}{S/n},
\]
take \(n=4\), \(S=20\). The vector \((5,5,5,5)\) has \(T=0\)
and inclusive upper tail one. The vector \((0,0,0,20)\) has \(T=60\)
and exact conditional upper tail \(4^{1-20}\).
\end{hbtheorem}

The two vectors have the same fitted rate and the same posterior distribution
of the rate, yet radically different allocations across cells. With
\(a=b=1\), both give the posterior \(\operatorname{Gamma}(21,5)\).
The conditional check deliberately removes information about the total:
it assesses equal-rate allocation, not the gamma prior or total-count fit.
A clustered vector may indicate rate heterogeneity or dependence; the tail
probability does not identify which mechanism caused it. This is an editorial
worked example supplementing the article's applications.

\section{The article's constrained-data applications}
\label{sec:p15-ppc-constrained-applications}

\subsection{Tomographic reconstruction}
In the motivating experiment, counts in detector bins are modeled as
independent \(Y_i\sim\operatorname{Poisson}(\theta_i)\), with
\(\theta=Ag+r\), known nonnegative detector operator \(A\) and
correction vector \(r\), and nonnegative image \(g\). A discretized
image has more pixel values than detector bins; nonnegativity can nevertheless
prevent exact interpolation of noisy counts. The article reports
\(n=22\,464\), about six million counts, and a fitted Pearson discrepancy
near \(30\,000\) \cite[sec.~1.2]{GelmanMengStern1996Main}.
For comparison, \(n+2\sqrt{2n}\) is about \(22\,888\).
The Poisson-to-normal reference here is an approximation justified by large
counts in almost all bins, not an exact theorem for arbitrary Poisson means.

The diagnostic problem persists when an improved fit gives a discrepancy
near \(22\,000\) or \(20\,000\): substantial fitting has taken place,
but the number of effective fitted directions is not obtained by counting
pixels. Possible failures include detector calibration, corrections,
dependence, excess Poisson variation and image discretization. A check alone
does not choose between them. The article uses this case as motivation; it
does not supply a posterior predictive simulation of this image. Nor does
nonnegativity alone imply that every large-exposure experiment eventually
interpolates exactly: that claim also needs the true mean to have an
appropriate interior position in the attainable mean set.

\subsection{Increasing and convex mortality rates}
For ages \(t=35,\ldots,64\), let \(N_t\) denote exposure and \(y_t\)
the number of deaths. The data used by the article are:
\begin{center}\small
\begin{tabular}{rrr@{\qquad}rrr}
\hline
Age & Exposure & Deaths & Age & Exposure & Deaths\\\hline
35&1771.5&3&50&1516.0&4\\
36&2126.5&1&51&1371.5&7\\
37&2743.5&3&52&1343.0&4\\
38&2766.0&2&53&1304.0&4\\
39&2463.0&2&54&1232.5&11\\
40&2368.0&4&55&1204.5&11\\
41&2310.0&4&56&1113.5&13\\
42&2306.5&7&57&1048.0&12\\
43&2059.5&5&58&1155.0&12\\
44&1917.0&2&59&1018.5&19\\
45&1931.0&8&60&945.0&12\\
46&1746.5&13&61&853.0&16\\
47&1580.0&8&62&750.0&12\\
48&1580.0&2&63&693.0&6\\
49&1467.5&7&64&594.0&10\\\hline
\end{tabular}
\end{center}
The half units reflect the source's treatment of partial-year exposure.
They are exposure measurements, not fractional binomial trial counts.
The article motivates a binomial model and then uses the Poisson
approximation \(Y_t\mid\theta_t\sim\operatorname{Poisson}(N_t\theta_t)\).
Write
\[
 C=\{\theta:0\leq\theta_{35}\leq\cdots\leq\theta_{64},\quad
 \theta_{t+1}-2\theta_t+\theta_{t-1}\geq0\ (36\leq t\leq63)\}.
\]
This is a rate model with increasing, discretely convex rates. An additional
upper bound of one specifies a different, bounded rate region if mortality
probabilities rather than the unbounded Poisson approximation are required.

\begin{hbtheorem}{thm:p15-ppc-mortality-posterior}{proved=annex:p15.ppc-mortality-posterior}{A proper constrained posterior from the rate likelihood}
For \(N_t>0\), nonnegative integer \(y_t\), and Lebesgue density
proportional to \(\mathbf1_C(\theta)\), the posterior density
\[
 \pi(\theta\mid y)=Z^{-1}\mathbf1_C(\theta)
      \prod_{t=35}^{64}\theta_t^{y_t}e^{-N_t\theta_t}
\]
has \(0<Z<\infty\). It is obtained by conditioning independent gamma
variables of shapes \(y_t+1\) and rates \(N_t\) on membership in \(C\).
For fixed \(N_t\) and a posterior draw \(\theta\), the source's
replication draws independent \(Y'_t\sim\operatorname{Poisson}(N_t\theta_t)\).
\end{hbtheorem}

The article samples the constrained posterior by random-walk Metropolis,
starting nine chains: three at the maximum likelihood fit and three at each
of two rough linear and quadratic starting curves. Its convergence assessment
compares between-chain and within-chain variation. This describes the
reported computation, not a guarantee from nine starts.

For this model use
\(X^2(y,\theta)=\sum_t(y_t-N_t\theta_t)^2/(N_t\theta_t)\).
Maximizing likelihood and minimizing Pearson discrepancy are distinct
optimizations. The following construction specifies both problems exactly.

\begin{hbtheorem}{thm:p15-ppc-mortality-fit}{proved=annex:p15.ppc-mortality-fit}{Convex computation of the two mortality fits}
For positive counts \(y_t\) and positive exposures \(N_t\), every
positive increasing convex rate vector on ages 35--64 has a unique
representation
\[
 \theta_t=\frac1{1000}\left[
 \eta_0+(t-35)\eta_1+
 \sum_{k=1}^{28}\eta_{k+1}(t-35-k)_+\right],
\]
where \((x)_+=\max(x,0)\), \(\eta_0>0\), and
\(\eta_1,\ldots,\eta_{29}\geq0\). Set
\(\mu_t(\eta)=N_t\theta_t\). The negative log likelihood, up to
data-only constants, and the Pearson objective are
\[
 L(\eta)=\sum_t[\mu_t(\eta)-y_t\log\mu_t(\eta)],\qquad
 P(\eta)=\sum_t\frac{(y_t-\mu_t(\eta))^2}{\mu_t(\eta)}.
\]
Each objective is strictly convex and attains a unique minimum on this
region. For either objective \(F\), a feasible vector is its minimum
exactly when
\[
 \frac{\partial F}{\partial\eta_j}=0\quad\hbox{if }\eta_j>0,
 \qquad
 \frac{\partial F}{\partial\eta_j}\geq0\quad\hbox{if }\eta_j=0.
\]
The derivatives are obtained from
\[
 \frac{\partial L}{\partial\eta_j}
 =\sum_t\left(1-\frac{y_t}{\mu_t}\right)
                   \frac{\partial\mu_t}{\partial\eta_j},\qquad
 \frac{\partial P}{\partial\eta_j}
 =\sum_t\left(1-\frac{y_t^2}{\mu_t^2}\right)
                   \frac{\partial\mu_t}{\partial\eta_j}.
\]
\end{hbtheorem}

For the printed data, numerical optimization of these exact Poisson
objectives gives Pearson discrepancies approximately \(30.68545\) at
the likelihood fit and \(29.34390\) at the Pearson fit. The latter
rounds to \(29.3\). These calculations use the nonnegative rate region
just specified; the fitted rates are also below one. Finite optimizer
tolerances do not turn a numerical answer into an exact real number.

The article's reported simulation reference calculations give:
\begin{center}\small
\begin{tabular}{lrrr}
\hline
Reference for minimum discrepancy & Mean & Variance & Upper tail\\\hline
Sampling at the fitted likelihood estimate &23.0&43.4&0.16\\
Posterior predictive, refitting every replicate&21.1&39.6&0.10\\\hline
\end{tabular}
\end{center}
The realized-discrepancy tail is reported as 0.063. These are the article's
simulation results, not exact distribution constants. The tentative
\(\chi^2_{26}\) suggested by four apparent fitted curve parameters would
instead have mean 26 and variance 52; the simulation directly challenges
that shortcut.
The likelihood fit lies on the boundary of the constraint region. The
source argues that interior parameter values tend to generate data fitting
the constraints more readily, and thus motivate a lower reference
distribution than sampling at the boundary fit. The reported change from
23.0 to 21.1 in its reference mean illustrates that reasoning in this case;
it is not a general stochastic-order theorem for arbitrary constraints.

The source's Figures 2--7 distinguish four objects: noisy observed rates,
the constrained likelihood fit, posterior latent-rate curves, and replicated
noisy rates. They are not interchangeable uncertainty bands. The replicated
rates tend to exceed the observations at later ages; a scalar tail alone
hides that localization. Screening high-risk older applicants and the
geometry induced by a uniform constrained prior are possible explanations.
The comparison of 0.10 and 0.063 suggests a role for the prior, but does
not prove that the likelihood is adequate or identify the defective assumption
\cite[secs.~3.1--3.2]{GelmanMengStern1996Main}.

\section{The latent-class application}
\label{sec:p15-ppc-latent-classes}

The article examines 93 infants, classified by motor activity and crying
at four months and fear of unfamiliar stimuli at fourteen months.
Let \(a\in\{1,2,3,4\}\), \(b,c\in\{1,2,3\}\) denote the three
categories. The observed \((y_{ab1},y_{ab2},y_{ab3})\) are:
\begin{center}\small
\begin{tabular}{rr|rrr@{\qquad}rr|rrr}
\hline
\(a\)&\(b\)&1&2&3&\(a\)&\(b\)&1&2&3\\\hline
1&1&5&4&1&3&1&3&3&4\\
1&2&0&1&2&3&2&0&2&3\\
1&3&2&0&2&3&3&1&1&7\\
2&1&15&4&2&4&1&2&1&2\\
2&2&2&3&1&4&2&0&1&3\\
2&3&4&4&2&4&3&0&3&3\\\hline
\end{tabular}
\end{center}
For an integer \(K\geq1\), a \(K\)-class model has weights \(w_h\geq0\), \(\sum_h w_h=1\),
and within-class probability vectors \(u_h,v_h,z_h\), of lengths
4, 3 and 3. Its cell probabilities are
\[
 p_{abc}(\theta)=\sum_{h=1}^K w_hu_{ha}v_{hb}z_{hc},\qquad
 Y\mid\theta\sim\operatorname{Multinomial}(93;p(\theta)).
\]
Conditional independence within each class does not imply independence in
the mixture. The source's two-class interpretation associates relatively
low scores with an uninhibited group and relatively high scores with an
inhibited group. These are model-based interpretations of the study's
class structure; a fit check does not establish biological kinds.

\begin{hbdefinition}{def:p15-ppc-dirichlet}{A Dirichlet probability vector}
For an integer \(m\geq2\) and positive parameters
\(\kappa_1,\ldots,\kappa_m\), put \(\kappa_+=\sum_i\kappa_i\).
The Dirichlet law on the probability simplex has density
\[
 \frac{\Gamma(\kappa_+)}{\prod_i\Gamma(\kappa_i)}
       \prod_{i=1}^m x_i^{\kappa_i-1},\qquad
 x_i>0,\quad \sum_i x_i=1,
\]
relative to Lebesgue measure in the first \(m-1\) coordinates,
with \(x_m=1-\sum_{i<m}x_i\). For \(m=1\), use the point mass
at \(x_1=1\). The normalized parameter vector specifies the center;
their sum is the concentration parameter. Each class-specific prior below
has its own explicitly chosen parameter vector.
\end{hbdefinition}

\begin{hbdefinition}{def:p15-ppc-multinomial-deviance}{Deviance and latent-class augmentation}
For an observed table with total \(N>0\), define
\[
 D(y,\theta)=2\sum_{a,b,c} y_{abc}
              \log\frac{y_{abc}}{Np_{abc}(\theta)}.
\]
A zero-count term is zero; a positive count with zero model probability
gives discrepancy \(+\infty\). The extended-valued convention can be
handled by its measurable comparison events. With positive Dirichlet
hyperparameters all cell probabilities are positive almost surely.
Introduce latent class counts \(n_{habc}\) summing over \(h\) to
\(y_{abc}\). Given \(y,\theta\), the counts in each cell have a
multinomial distribution with probabilities proportional to
\(w_hu_{ha}v_{hb}z_{hc}\).
\end{hbdefinition}

\begin{hbtheorem}{thm:p15-ppc-mixture}{proved=annex:p15.ppc-mixture}{The augmented posterior and the discrepancy calculation}
Give \(w\) and each of \(u_h,v_h,z_h\) independent proper Dirichlet
priors with strictly positive parameter vectors. Conditional on the latent
counts, their posterior distributions are independent Dirichlet laws with
each prior parameter increased by its corresponding count margin.
The displayed deviance is twice the log likelihood ratio of the saturated
multinomial fit to \(p(\theta)\), is nonnegative, and is invariant under
a simultaneous permutation of all class labels. A replication with the
same latent-class parameters and sample size has exactly the multinomial
law displayed above; it does not fix the realized latent class counts.
\end{hbtheorem}

For \(K=1,2,3,4\), the article reports minimized deviances
48.761, 14.150, 9.109 and 4.718. Its parameter-count residual dimensions
are 28, 20, 12 and 4. These dimension counts do not prove corresponding
chi-square reference distributions: absent classes, coincident classes,
boundary parameters and nonunique representations prevent casual use of
regular likelihood-ratio asymptotics.

The fitted two-class model is checked using independent Dirichlet prior
blocks centered on the psychological predictions, with weak concentration.
Their numerical specification is given in the separate Rubin--Stern
1994 analysis referred to by the main article.
Posterior computation uses latent-class data augmentation, ten dispersed
starts and a multiple-chain convergence assessment. The probabilities below
are results reported for that specification; the generic positive-Dirichlet
construction above does not determine those numerical results.

The reported predictive reference for \(D_{\min}\) has mean 23.4,
variance 45.3 and upper tail 0.93; the realized-discrepancy tail is 0.74.
Figures 8 and 9 display, respectively, the reference histogram with the
observed minimum marked and the paired-discrepancy scatter with its diagonal.
The numerical contrast with \(\chi^2_{20}\), whose mean and variance
would be 20 and 40, illustrates why a fitted dimension count is insufficient.
The 0.74 realized-discrepancy calculation is reported with 500 draws.

For the one-class independence model the corresponding tails are 0.024
and 0.058. Treating the conventional 0.05 threshold as a sharp substantive
divide would conceal their common indication of poor fit. The two-class
checks find no comparable discrepancy along the checked directions; they
neither establish the model as true nor select the number of classes by
posterior model probability. The source also reports two likelihood modes
for the two-class fit. A calculation trapped at one mode cannot claim to
average the full posterior \cite[secs.~3.3--3.4]{GelmanMengStern1996Main}.

\section{Prior, posterior and mixed replications}
\label{sec:p15-ppc-replication-choices}

For fixed design information, the three reference experiments have different
shared quantities:
\begin{center}\small
\begin{tabular}{p{.15\linewidth}p{.40\linewidth}p{.30\linewidth}}
\hline
Replication & Shared with the observed experiment & Newly drawn\\\hline
Posterior & \(\theta\), uncertain given \(y\)&data given \(\theta\)\\
Prior & model and proper prior&\(\theta'\sim\Pi\), then data\\
Mixed & hyperparameter \(\alpha\), uncertain given \(y\)&group parameter, then data\\\hline
\end{tabular}
\end{center}
For a hierarchy with hyperprior \(\Lambda(d\alpha)\), parameter kernel
\(K(\alpha,d\theta)\), and data kernel \(P_\theta\), let
\(\Lambda_y\) be the posterior law of \(\alpha\). The mixed
predictive data law with fixed design is
\[
 M^{\rm mix}_y(B)=\int\!\int P_{\theta'}(B)
                 K(\alpha,d\theta')\Lambda_y(d\alpha).
\]
It checks a new group from the same inferred population distribution,
whereas the posterior replication checks a repeat at the same latent group
parameters. A realization-dependent discrepancy in a mixed experiment must
specify whether its observed and replicated arguments use \(\theta\),
\(\theta'\), or \(\alpha\); the posterior shared-parameter formula
does not decide that question.

If \(A\) is random rather than fixed by design, specify the full joint
replication model and then condition it on the required equality of auxiliary
values. For example, a fresh prior-drawn parameter, conditioned on
\(A(Y')=a\), has weights proportional to
\(\Pi(d\theta')P_{\theta'}(A(Y')=a)\) in the discrete positive-mass
case. Drawing from the raw prior and merely conditioning the data kernel
need not be that same experiment. This explicit choice prevents an auxiliary
conditioning convention from silently changing the meaning of a check.

\subsection{A diffuse prior and a maximum statistic}
\begin{hbtheorem}{thm:p15-ppc-normal-maximum}{proved=annex:p15.ppc-normal-maximum}{Exact formulas for the normal maximum example}
Let \(n\geq1\) be an integer, let \(Y_i\mid\theta\), \(1\leq i\leq n\), be independent \(N(\theta,1)\), let
\(\theta\) have the uniform prior on \([-L,L]\), \(L>0\), and set
\(T(y)=\max_i|y_i|\). The posterior is
\(N(\bar y,1/n)\) restricted and normalized on \([-L,L]\).
For \(t\geq0\), write
\[
 G_t(\theta)=\{\Phi(t-\theta)-\Phi(-t-\theta)\}^n.
\]
The predictive lower-tail probabilities of \(T(Y')\) are
\[
 F_{\rm post}(t)=\int G_t(\theta)\Pi_y(d\theta),\qquad
 F_{\rm prior}(t)=\frac1{2L}\int_{-L}^L G_t(\theta)\,d\theta.
\]
For fixed finite \(t\), \(F_{\rm prior}(t)\to0\) as \(L\to\infty\).
For fixed observed data the posterior converges to
\(N(\bar y,1/n)\), and its predictive probabilities converge to those
obtained from that normal posterior.
\end{hbtheorem}

The article uses \(n=100\), \(\bar y=5.1\), \(T(y)=8.1\) and
\(L=10^5\). Its effectively untruncated posterior simulation yields
149 upper-tail exceedances among 1,000 replicates. The prior replication
instead regards such a small absolute maximum as highly unusual because
the prior spreads means over a vast interval. These answer different
questions about what is repeated. The source also suggests the more nearly
pivotal statistic \(\max_i|y_i|-|\bar y|\); it does not claim that this
absolute-value adjustment is an exact pivot for every mean.
Numerically integrating the displayed normal-posterior reference gives
an upper tail of approximately \(0.13129403\). The reported simulation
fraction \(0.149\) is about \(1.66\) independent-draw Monte Carlo
standard deviations above that value. The integration and the historical
simulation fraction are different numerical approximations, rather than
two exact definitions of the reference probability.

\subsection{Prior sensitivity, relevance and decisions}
In the mixture example the article varies both the centers and concentration
of the Dirichlet priors. Weak alternatives give similar checks; a strongly
informative prior forcing inappropriate class profiles can make the realized
tail essentially zero. A trusted informative prior can sharpen a check of
the sampling model, while an incorrect strong prior can itself cause poor
prediction. Differences between minimum and realized checks are clues,
not a general identification theorem for which part of the model failed.

A discrepancy should bear on the intended inference. A model can predict
a median adequately while giving poor estimates of a total; checks of a
total and of a median then serve different uses. Displays of residuals,
tails, subgroup behavior and dependence complement a global discrepancy.
Sensitivity of conclusions to priors and likelihood choices, and substantive
plausibility of posterior inferences, remain additional checks beyond
posterior predictive agreement.

The source advocates examining many discrepancies descriptively without
treating their individual tails as a family of accept/reject decisions.
That practice supplies no automatic familywise-error guarantee. If selected
checks are used to claim controlled error, to choose one unusually small
tail after searching, or to make a decision with specified costs, the
selection mechanism and decision rule require their own analysis. Neither
a replicated-data plot nor a predictive tail supplies those missing claims.
Defects may be identified even when no adequate replacement model is yet
available; use of a defective approximation should retain the limits that
matter to its intended inference \cite[sec.~4]{GelmanMengStern1996Main}.
